\section{Sensitivity to the Modeling Assumptions}
\label{sec:robustness}
\label{sec:powercurve}
\label{sec:gaussian}
%% generated by mpce_results.py -- do not edit by hand
\begin{table}[!t]
\centering
\caption{Robustness to the Benchmark Model: Final Layouts Re-Evaluated, Not Re-Optimized}
\label{tab:robust-final}
\scriptsize\setlength{\tabcolsep}{2.4pt}
\begin{tabular}{lcccccc}
\toprule
& \multicolumn{2}{c}{$\Delta$ (\%)} & & & & \\
\cmidrule(lr){2-3}
Model & DS I & DS II & Rank & Best & $\bar\tau$ & Same \\
\midrule
Benchmark & +0.0 & +0.0 & 1.82 & PSO-VNS & -- & -- \\
Cubic curve & -17.7 & -36.9 & 1.82 & PSO-VNS & 0.95 & 94 \\
Cubic, 25 m/s cut-out & -21.5 & -37.0 & 1.80 & PSO-VNS & 0.95 & 93 \\
Gaussian wake & -0.3 & -0.4 & 1.97 & PSO & 0.66 & 54 \\
\bottomrule
\end{tabular}
\par\vspace{2pt}\parbox{\columnwidth}{\scriptsize All feasible final layouts of the 68 cases. $\Delta$: mean relative change of the objective (\%); Rank: average rank of PSO-VNS; Best: best-ranked method; $\bar\tau$: mean Kendall correlation between the benchmark and the alternative ordering within a case; Same: cases (\%) with unchanged best method. Ranks of all methods: supplementary material.}
\end{table}


\emph{Power curve and wake model:} We re-evaluated every feasible final layout of the main comparison, without re-optimization, with 1) a cubic power curve~\cite{Carrillo2013}, $f_{\rm c}(s)=P_{\text{rated}}(s^3-s_{\text{cut-in}}^3)/(s_{\text{rated}}^3-s_{\text{cut-in}}^3)$ between the cut-in and rated speeds, 2) the same curve with a 25~m/s cut-out, and 3) the Gaussian wake model of~\cite{Bastankhah2014} with growth rate $K_{\rm G}=0.04$ and the superposition and Weibull scaling of the benchmark (Section~\ref{S-sec:S-gauss} of the supplementary material). The absolute energy depends strongly on the model (Table~\ref{tab:robust-final}): the cubic curve gives \NCubicOverI\% (Data Set~I) and \NCubicOverII\% (Data Set~II) less expected power than the linear benchmark curve, so our energy values are benchmark values, not AEP predictions for a specific turbine. The ranking of the methods is stable under the cubic curves ($\bar\tau=\NTauCubic$; same best method in \NSameCubic\% of the cases), and PSO-VNS keeps the best average rank. It is less stable under the Gaussian model ($\bar\tau=\NTauGauss$; same best method in \NSameGauss\% of the cases), since layouts tuned to the sharp Jensen wake cone are judged differently by a smooth deficit: \NBestGauss{} ranks first (\NRankGaussPSO) and PSO-VNS \NPosGaussPSOVNS{} (\NRankGaussPSOVNS), both ahead of all other methods. The small advantage of the VNS phase over PSO is thus specific to the wake model it optimizes, which is a limit of PSO-VNS.
% CHECK-FINAL [C21]: robustness.Cubic/CubicCutout: tau >= 0.9, PSO-VNS rank position 1
% CHECK-FINAL [C22]: robustness.Gauss: tau < 0.8, same_best < 75, PSOC first, PSOBV second

\label{sec:spacing}
\label{sec:boundary}
\emph{Spacing and boundary handling:} The optimized layouts often lie on the $4D$ constraint inherited from the benchmark: of the feasible final layouts of the eight methods, \NSpFiveAll\% also satisfy $5D$ and \NSpSixAll\% satisfy $6D$ (\NSpFiveLarge\% and \NSpSixLarge\% for $N=11$--15), so they cannot be transferred to sites with a larger spacing without re-optimization.
% CHECK-FINAL [C23]: spacing_share.n11_15.pct_5d < 10 ("often lie on the 4D constraint")
At $5D$ and $6D$, a multi-start packing search fits fewer turbines into each circle, and the largest 500-m cases may be infeasible (constructive lower bounds in Table~\ref{S-tab:capacity}; a re-optimization of six cases with SSA and LX-SSA is given in Table~\ref{S-tab:spacing-authors}). The boundary rule also matters: an otherwise identical SSA that projects every turbine radially onto the circle instead of clipping the coordinates attains higher mean benchmark objectives than the original SSA in all six representative cases at 3,030 evaluations (Mann--Whitney test; \TBD{macro: range of mean differences and max.\ $p$ of the radial-projection vs.\ clipping SSA runs}). We therefore apply the rule of Section~\ref{sec:constraints} to all methods.
